\documentclass[12pt]{article}

\usepackage{sbc-template}
\usepackage[brazil]{babel}
\usepackage[utf8]{inputenc}
\usepackage{graphicx}
\usepackage{booktabs}
\usepackage{multirow}
\usepackage{url}
\usepackage{hyperref}
\usepackage{float}
\usepackage{xcolor}
\usepackage{tabularx}
\usepackage{amsmath}

\sloppy

\title{Análise de Vulnerabilidades de Segurança em Código Gerado por Modelos de Linguagem: Um Estudo Comparativo entre GPT e Claude em Prompts em Português e Inglês}

\author{Tonil A. Viera\inst{1}}

\address{Instituto de Ciências Exatas e Informática\\
  Pontifícia Universidade Católica de Minas Gerais (PUC Minas)\\
  Belo Horizonte, MG -- Brasil
  \email{tonil.viera@sga.pucminas.br}
}

\begin{document}

\maketitle

\begin{abstract}
This paper presents a comparative analysis of security vulnerabilities introduced by Large Language Models (LLMs) when generating source code from natural language prompts written in Portuguese (PT-BR) and English (EN-US). Two state-of-the-art models — GPT-5.4 Codex and Claude Sonnet 4.6 — were evaluated across 24 implementations spanning three programming languages (C++, Go, Python) and two application scenarios (Lead Management REST API and Music Streaming Service). Each implementation was subjected to automated static analysis using Cppcheck, Gosec, and Bandit, as well as a manual security code review following CWE classification. A total of 392 vulnerabilities were identified: 193 in Portuguese-prompted code and 199 in English-prompted code (3.1\% difference). Claude generated an average of 17.2 vulnerabilities per project versus 15.5 for GPT. CWE-798 (Hardcoded Credentials) was the only vulnerability present in all 24 implementations, confirming it as a systemic pattern in LLM-generated code regardless of language or model. The results suggest that prompt language influences complexity of generated code — particularly in C++ (more issues in PT) and Python (more issues in EN) — with linguistic descriptiveness and technical register as contributing factors.
\end{abstract}

\begin{resumo}
Este artigo apresenta uma análise comparativa de vulnerabilidades de segurança introduzidas por Modelos de Linguagem de Grande Escala (LLMs) ao gerar código-fonte a partir de prompts em Português (PT-BR) e Inglês (EN-US). Dois modelos de ponta — GPT-5.4 Codex e Claude Sonnet 4.6 — foram avaliados em 24 implementações abrangendo três linguagens de programação (C++, Go, Python) e dois cenários de aplicação (API REST de Gerenciamento de Leads e Serviço de Streaming de Música). Cada implementação foi submetida à análise estática automatizada com Cppcheck, Gosec e Bandit, além de revisão manual de segurança seguindo a classificação CWE. Foram identificadas 392 vulnerabilidades: 193 em código gerado por prompts em português e 199 em inglês (diferença de 3,1\%). O Claude gerou em média 17,2 vulnerabilidades por projeto contra 15,5 do GPT. A CWE-798 (Credenciais Codificadas) foi a única vulnerabilidade presente nas 24 implementações, confirmando-se como padrão sistêmico do código gerado por LLMs, independente de idioma ou modelo. Os resultados sugerem que o idioma do prompt influencia a complexidade do código gerado — especialmente em C++ (mais problemas em PT) e Python (mais problemas em EN) — com a descrição linguística e o registro técnico como fatores contribuintes.
\end{resumo}

%-------------------------------------------------------------------
\section{Introdução}
%-------------------------------------------------------------------

A adoção de Modelos de Linguagem de Grande Escala (LLMs) para geração de código tornou-se amplamente difundida entre desenvolvedores, equipes ágeis e estudantes de computação. Ferramentas como GitHub Copilot, ChatGPT e Claude Sonnet têm demonstrado capacidade de produzir implementações funcionais a partir de descrições em linguagem natural, acelerando significativamente o ciclo de desenvolvimento de software~\cite{pearce2022asleep}.

Entretanto, a qualidade de segurança do código gerado automaticamente por LLMs ainda é objeto de investigação. Estudos recentes indicam que esses modelos tendem a introduzir vulnerabilidades conhecidas — como credenciais codificadas em código-fonte, transmissão em texto claro, ausência de limitação de taxa e configurações permissivas de CORS — mesmo em implementações funcionalmente corretas~\cite{pearce2022asleep,khoury2023secure}. Essa tendência é especialmente relevante no contexto brasileiro, onde uma parcela significativa dos desenvolvedores utiliza o português como idioma nativo para redigir especificações e prompts.

A questão central deste trabalho é: \textit{a escolha do idioma do prompt — português (PT-BR) versus inglês (EN-US) — influencia o perfil de vulnerabilidades de segurança no código gerado por LLMs?} Para responder a essa questão, este estudo propõe um protocolo experimental controlado que gera 24 implementações completas (dois modelos $\times$ dois idiomas $\times$ dois cenários $\times$ três linguagens de programação), avalia cada implementação por análise estática automatizada e por revisão manual de segurança baseada em CWE, e compara os resultados entre idiomas, modelos e linguagens de programação.

As principais contribuições deste trabalho são:
\begin{itemize}
  \item Um protocolo reproduzível para avaliação de segurança de código gerado por LLMs;
  \item Dados empíricos sobre a influência do idioma do prompt no perfil de vulnerabilidades;
  \item Análise qualitativa dos padrões de vulnerabilidade mais recorrentes em código LLM;
  \item Uma hipótese linguística para explicar as diferenças observadas entre PT-BR e EN-US.
\end{itemize}

O restante do artigo está organizado da seguinte forma: a Seção~\ref{sec:trabalhos} apresenta trabalhos relacionados; a Seção~\ref{sec:metodologia} descreve a metodologia; a Seção~\ref{sec:implementacoes} documenta as implementações; a Seção~\ref{sec:estatica} apresenta os resultados da análise estática; a Seção~\ref{sec:manual} descreve a revisão manual; a Seção~\ref{sec:resultados} analisa os resultados; e as Seções~\ref{sec:conclusao} e~\ref{sec:futuros} apresentam conclusão e trabalhos futuros.

%-------------------------------------------------------------------
\section{Trabalhos Relacionados}
\label{sec:trabalhos}
%-------------------------------------------------------------------

Pearce et al.~\cite{pearce2022asleep} realizaram o primeiro estudo sistemático sobre vulnerabilidades introduzidas por GitHub Copilot, analisando 89 cenários de geração de código e identificando que cerca de 40\% das sugestões continham vulnerabilidades de segurança. Os autores concluem que o Copilot pode ser considerado "dormindo no volante" — funcionalmente produtivo, mas inseguro.

Khoury et al.~\cite{khoury2023secure} expandiram esse estudo para ChatGPT (GPT-3.5), avaliando 21 programas em cinco linguagens. Todos os programas gerados apresentaram ao menos uma vulnerabilidade de segurança, com as categorias CWE-20 (Validação de Entrada), CWE-787 (Escrita Fora dos Limites) e CWE-79 (XSS) sendo as mais recorrentes. Os autores destacam que o modelo demonstra conhecimento teórico de segurança, mas não o aplica consistentemente na geração de código.

Tony et al.~\cite{tony2023llmseceval} desenvolveram o LLMSecEval, um benchmark de 150 cenários de geração de código para avaliação de segurança de LLMs, cobrindo os Top-25 CWEs do MITRE. Os resultados mostram que modelos maiores tendem a gerar código mais seguro, mas nenhum modelo testado está isento de vulnerabilidades críticas.

No contexto de segurança em API REST, Yao et al.~\cite{yao2024survey} apresentam uma revisão sistemática sobre código LLM-gerado e concluem que padrões de autenticação (JWT, OAuth) são frequentemente implementados de forma insegura, com destaque para a ausência de verificação de expiração de tokens e o uso de segredos previsíveis.

Diferentemente dos trabalhos citados, este estudo contribui com a variável do \textit{idioma do prompt} como fator de análise, preenche a lacuna de avaliações em PT-BR e aplica tanto análise estática automatizada quanto revisão manual alinhada ao CWE, cobrindo dois cenários de aplicação completos em três linguagens de programação.

%-------------------------------------------------------------------
\section{Metodologia}
\label{sec:metodologia}
%-------------------------------------------------------------------

\subsection{Modelos Avaliados}

Foram selecionados dois modelos de linguagem de ponta:
\begin{itemize}
  \item \textbf{GPT-5.4 Codex (OpenAI)}: modelo de alta capacidade com especialização em geração de código, acessado via interface de chat padrão;
  \item \textbf{Claude Sonnet 4.6 (Anthropic)}: modelo conversacional com forte desempenho em tarefas de programação, acessado via interface de chat e Claude Code CLI.
\end{itemize}

\subsection{Cenários de Aplicação}

Dois cenários de aplicação foram definidos para garantir cobertura de padrões de segurança distintos:

\textbf{Cenário 1 — Gerenciamento de Leads (Leads):} API REST para gerenciamento de contatos/leads comerciais, com autenticação JWT, operações CRUD, importação/exportação CSV/Excel e integração com Firebase Firestore como banco de dados.

\textbf{Cenário 2 — Streaming de Música (Stream):} Serviço de streaming de áudio com separação de papéis (usuário/artista), upload de arquivos MP3, reprodução em tempo real e autenticação baseada em tokens.

Ambos os cenários envolvem mecanismos de autenticação, manipulação de arquivos, gerenciamento de sessão e exposição de dados sensíveis — áreas que historicamente concentram vulnerabilidades de segurança.

\subsection{Linguagens de Programação}

Três linguagens foram selecionadas por sua heterogeneidade em paradigmas de segurança e popularidade no mercado:
\begin{itemize}
  \item \textbf{C++}: linguagem de sistema, gerenciamento manual de memória, maior superfície para vulnerabilidades de baixo nível;
  \item \textbf{Go}: linguagem moderna com garbage collector, tipagem estática e bibliotecas padrão de criptografia;
  \item \textbf{Python}: linguagem dinâmica com vasto ecossistema de frameworks web (FastAPI, Flask).
\end{itemize}

\subsection{Design Experimental}

O experimento segue um design fatorial completo $2 \times 2 \times 3$ (modelos $\times$ idiomas $\times$ linguagens), resultando em 12 configurações por cenário e 24 implementações no total (Tabela~\ref{tab:design}).

\begin{table}[H]
\centering
\caption{Design experimental: 24 implementações}
\label{tab:design}
\begin{tabular}{llllll}
\toprule
\textbf{Cenário} & \textbf{Linguagem} & \textbf{Idioma PT} & & \textbf{Idioma EN} & \\
 & & Claude & GPT & Claude & GPT \\
\midrule
Leads & C++ & \checkmark & \checkmark & \checkmark & \checkmark \\
Leads & Go & \checkmark & \checkmark & \checkmark & \checkmark \\
Leads & Python & \checkmark & \checkmark & \checkmark & \checkmark \\
Stream & C++ & \checkmark & \checkmark & \checkmark & \checkmark \\
Stream & Go & \checkmark & \checkmark & \checkmark & \checkmark \\
Stream & Python & \checkmark & \checkmark & \checkmark & \checkmark \\
\bottomrule
\end{tabular}
\end{table}

\subsection{Protocolo de Geração de Código}

Para cada uma das 24 células do design, foi elaborado um prompt equivalente nos dois idiomas descrevendo os requisitos funcionais do cenário (autenticação, CRUD, upload/export, integração com Firebase). Os prompts em PT-BR e EN-US descrevem as mesmas funcionalidades, diferindo apenas no idioma. O código gerado foi executado localmente para verificação de compilação/execução antes da análise de segurança.

\subsection{Instrumentos de Análise}

\textbf{Análise Estática Automatizada:} Três ferramentas foram utilizadas conforme a linguagem alvo:
\begin{itemize}
  \item \textbf{Cppcheck 2.x} para C++: análise de memória, ponteiros nulos, buffer overflow e uso de APIs inseguras;
  \item \textbf{Gosec v2.x} para Go: análise de criptografia fraca, credenciais hardcoded, injeção em logs (CWE-117) e erros não tratados;
  \item \textbf{Bandit 1.x} para Python: análise de uso de funções inseguras, criptografia inadequada e configurações de rede permissivas.
\end{itemize}

\textbf{Revisão Manual de Segurança:} Para cada implementação, foi conduzida uma revisão manual sistemática do código-fonte, mapeando cada vulnerabilidade encontrada ao identificador CWE correspondente, classificando-a por severidade (Crítica, Alta, Média, Baixa) e documentando o vetor de exploração. A revisão cobriu exclusivamente o código da aplicação, excluindo bibliotecas de terceiros.

%-------------------------------------------------------------------
\section{Implementações}
\label{sec:implementacoes}
%-------------------------------------------------------------------

As 24 implementações foram geradas entre março e abril de 2026. Cada implementação resultou em um projeto completo com backend funcional, incluindo configuração de build (CMakeLists.txt para C++, go.mod para Go, requirements.txt para Python), código de integração com Firebase Firestore e, onde aplicável, interfaces de usuário em HTML/CSS/JavaScript.

Os projetos foram organizados na estrutura de diretórios:
\begin{verbatim}
tcc/
  Portugues/{Leads,Stream}/{C++,Go,Python}/{Claude,GPT}/
  English/{Leads,Stream}/{C++,Go,Python}/{Claude,GPT}/
\end{verbatim}

\textbf{Cenário Leads:} As implementações em C++ utilizaram cpp-httplib ou Crow como framework HTTP. As em Go utilizaram chi, Gin ou servidor HTTP nativo. As em Python utilizaram FastAPI ou Flask. Todas implementam os endpoints: \texttt{POST /api/auth/register}, \texttt{POST /api/auth/login}, \texttt{GET/POST/PUT/DELETE /api/leads}, \texttt{POST /api/leads/import} e \texttt{GET /api/leads/export}.

\textbf{Cenário Stream:} As implementações em C++ utilizaram cpp-httplib. As em Go utilizaram chi ou servidor HTTP nativo com SQLite ou JSON para catálogo. As em Python utilizaram FastAPI com SQLite ou Firebase. Todos os projetos implementam: \texttt{POST /api/auth/register} e \texttt{/login}, \texttt{POST /api/artist/songs} (upload), \texttt{GET /api/songs} (listagem), \texttt{GET /api/songs/{id}/stream} (streaming com Range headers).

A integração com Firebase Firestore foi comum a todas as implementações do cenário Leads. No cenário Stream, alguns modelos optaram por SQLite local combinado com sistema de arquivos, enquanto outros mantiveram Firebase para armazenamento de metadados.

%-------------------------------------------------------------------
\section{Análise Estática Automatizada}
\label{sec:estatica}
%-------------------------------------------------------------------

A análise estática foi executada sobre o código-fonte de cada projeto utilizando os resultados armazenados nos arquivos \texttt{report.txt} de cada implementação. Os resultados são apresentados na Tabela~\ref{tab:estatica}.

\begin{table}[H]
\centering
\caption{Resultados da análise estática por ferramenta e idioma}
\label{tab:estatica}
\begin{tabular}{llrr}
\toprule
\textbf{Linguagem} & \textbf{Ferramenta} & \textbf{Total PT} & \textbf{Total EN} \\
\midrule
C++ & Cppcheck & 8 & 9 \\
Go & Gosec & 31 & 45 \\
Python & Bandit & 13 & 15 \\
\midrule
\textbf{Total} & & \textbf{52} & \textbf{69} \\
\bottomrule
\end{tabular}
\end{table}

O Gosec foi a ferramenta com maior número de ocorrências reportadas, principalmente devido ao flagging agressivo de injeção em logs (CWE-117) e erros não tratados em chamadas de API (CWE-703) em projetos Go. Esses tipos de alertas possuem alto índice de falsos positivos em código Go convencional, pois o Gosec aplica análise de fluxo de dados que frequentemente captura padrões válidos.

Os projetos C++ em português apresentaram ligeiramente menos alertas de Cppcheck (8 vs 9 em inglês), em parte porque a geração GPT/PT do cenário Leads/C++ utilizou uma biblioteca externa (\texttt{json.hpp}) que foi erroneamente incluída no escopo de análise do Cppcheck, sendo posteriormente descartada da contagem por não fazer parte do código da aplicação.

A análise estática capturou um subconjunto restrito das vulnerabilidades totais identificadas na revisão manual, confirmando que ferramentas automatizadas são complementares, e não substitutivas, à revisão humana de segurança.

%-------------------------------------------------------------------
\section{Revisão Manual de Segurança}
\label{sec:manual}
%-------------------------------------------------------------------

A revisão manual de segurança foi conduzida de forma independente para cada uma das 24 implementações, sem consulta prévia aos relatórios de análise estática, para evitar viés de confirmação. O processo seguiu o seguinte protocolo:

\begin{enumerate}
  \item Leitura completa do código-fonte da aplicação (excluindo bibliotecas de terceiros);
  \item Identificação de cada padrão de vulnerabilidade com mapeamento ao CWE correspondente;
  \item Classificação de severidade segundo o CVSS 3.1 (Crítica/Alta/Média/Baixa);
  \item Documentação de vetor de exploração para cada vulnerabilidade identificada;
  \item Registro de observações positivas (boas práticas encontradas).
\end{enumerate}

As revisões foram documentadas nos arquivos \texttt{security\_review.txt} de cada projeto. O total de vulnerabilidades identificadas manualmente foi de \textbf{392}, sendo 193 em projetos PT-BR e 199 em projetos EN-US.

As categorias CWE mais recorrentes na revisão manual incluem: CWE-798 (Credenciais Hardcoded), CWE-319 (Transmissão em Texto Claro), CWE-307/770 (Ausência de Rate Limiting), CWE-942 (CORS Permissivo), CWE-352 (CSRF), CWE-521 (Política de Senha Fraca) e CWE-613 (Expiração de Sessão Insuficiente).

%-------------------------------------------------------------------
\section{Resultados e Discussão}
\label{sec:resultados}
%-------------------------------------------------------------------

\subsection{Visão Geral: Contagem por Projeto}

A Tabela~\ref{tab:resultados_completos} apresenta o número de vulnerabilidades identificadas na revisão manual para cada um dos 24 projetos, organizados por cenário, linguagem de programação, idioma do prompt e modelo de linguagem utilizado.

\begin{table}[H]
\centering
\caption{Vulnerabilidades por projeto — Revisão Manual (CWE)}
\label{tab:resultados_completos}
\begin{tabular}{llrrrr}
\toprule
\multirow{2}{*}{\textbf{Cenário}} & \multirow{2}{*}{\textbf{Ling.}} &
\multicolumn{2}{c}{\textbf{PT-BR}} & \multicolumn{2}{c}{\textbf{EN-US}} \\
 & & \textbf{Claude} & \textbf{GPT} & \textbf{Claude} & \textbf{GPT} \\
\midrule
Leads & C++ & 27 & 18 & 16 & 14 \\
Leads & Go  & 14 & 14 & 19 & 12 \\
Leads & Python & 11 & 14 & 23 & 13 \\
\midrule
Stream & C++ & 18 & 17 & 18 & 25 \\
Stream & Go  & 16 & 14 & 14 & 17 \\
Stream & Python & 16 & 14 & 14 & 14 \\
\midrule
\textbf{Total PT} & & \textbf{102} & \textbf{91} & & \\
\textbf{Total EN} & & & & \textbf{104} & \textbf{95} \\
\textbf{Total geral} & & \multicolumn{2}{c}{\textbf{193}} & \multicolumn{2}{c}{\textbf{199}} \\
\bottomrule
\end{tabular}
\end{table}

\subsection{Médias por Agrupamento}

A Tabela~\ref{tab:medias} consolida as médias de vulnerabilidades por projeto segundo diferentes agrupamentos analíticos.

\begin{table}[H]
\centering
\caption{Média de vulnerabilidades por projeto}
\label{tab:medias}
\begin{tabular}{llrr}
\toprule
\textbf{Dimensão} & \textbf{Grupo} & \textbf{Total} & \textbf{Média/projeto} \\
\midrule
Idioma & PT-BR & 193 & 16,1 \\
Idioma & EN-US & 199 & 16,6 \\
\midrule
Modelo & Claude & 206 & 17,2 \\
Modelo & GPT    & 186 & 15,5 \\
\midrule
Linguagem & C++    & 153 & 19,1 \\
Linguagem & Go     & 120 & 15,0 \\
Linguagem & Python & 119 & 14,9 \\
\midrule
Cenário & Leads  & 195 & 16,3 \\
Cenário & Stream & 197 & 16,4 \\
\bottomrule
\end{tabular}
\end{table}

A diferença entre os idiomas é de apenas 3,1\% (199 vs 193), estatisticamente pequena. A diferença entre modelos é mais expressiva: Claude gerou em média 17,2 vulnerabilidades por projeto contra 15,5 do GPT (diferença de 11\%). C++ se destacou como a linguagem com maior média (19,1), seguida de Go (15,0) e Python (14,9). Os dois cenários apresentaram médias praticamente idênticas (16,3 vs 16,4), indicando que a complexidade de segurança dos dois domínios é comparável.

\subsection{Vulnerabilidades Mais Críticas}

\subsubsection{CWE-798: Credenciais Hardcoded (Universal)}

A vulnerabilidade CWE-798 — uso de credenciais codificadas diretamente no código-fonte — foi identificada em \textbf{todas as 24 implementações} analisadas (100\% de ocorrência). Em todos os casos, dois padrões se repetiram:

\begin{enumerate}
  \item \textbf{Segredo JWT hardcoded}: todos os projetos incluíam um valor padrão previsível para o segredo de assinatura de tokens JWT (ex.: \texttt{"lead-manager-secret-key-change-in-production"}, \texttt{"change-this-secret-in-production"}, \texttt{"troque-esta-chave-em-producao"}). Esse padrão permite que um atacante forje tokens JWT arbitrários e assuma qualquer conta no sistema.
  \item \textbf{Caminho absoluto para credenciais Firebase}: todos os projetos incluíram, como valor padrão, o caminho absoluto para o arquivo de credenciais do Firebase Admin SDK (\texttt{C:/Users/tonil/Desktop/tcc/lead-manager-54595-firebase-adminsdk-fbsvc-375221693a.json}), expondo o nome do projeto Firebase e a estrutura de diretórios do desenvolvedor.
\end{enumerate}

A universalidade da CWE-798 não é acidental: reflete um padrão comportamental dos LLMs de priorizar conveniência de desenvolvimento (código que "funciona imediatamente") sobre segurança, inserindo placeholders facilmente descobertos.

\subsubsection{CWE-319: Transmissão em Texto Claro}

CWE-319 — ausência de TLS/HTTPS — foi identificada em 20 dos 24 projetos (83\%), afetando principalmente implementações em C++ e Go, onde o desenvolvedor precisa configurar explicitamente um servidor SSL. Em C++, todos os oito projetos utilizaram \texttt{httplib::Server} em vez de \texttt{httplib::SSLServer}. Em Go, todos os oito projetos utilizaram \texttt{http.ListenAndServe} sem TLS. Nos projetos Python, a responsabilidade pelo TLS é geralmente delegada a um proxy reverso (nginx), o que pode explicar a menor incidência nessa linguagem.

\subsubsection{CWE-307/770: Ausência de Rate Limiting}

A ausência de limitação de taxa nos endpoints de autenticação (\texttt{/login}, \texttt{/register}) foi identificada em 22 dos 24 projetos (92\%). Esse padrão viabiliza ataques de força bruta e credential stuffing. Os dois projetos que implementaram alguma forma de proteção o fizeram parcialmente (limitação de tamanho de payload, sem limitação de frequência).

\subsubsection{CWE-942: CORS Permissivo}

Configurações de CORS excessivamente permissivas (\texttt{allow\_origins=["*"]} combinado com \texttt{allow\_credentials=True}) foram encontradas em 20 dos 24 projetos (83\%). Essa combinação é especialmente perigosa pois viola a especificação CORS (que proíbe credenciais com origem curinga) e pode ser explorada dependendo da implementação do servidor.

\subsection{Padrões por Idioma do Prompt}

\subsubsection{Diferenças por Linguagem de Programação}

A Tabela~\ref{tab:por_linguagem} apresenta os totais de vulnerabilidades desagregados por linguagem de programação e idioma do prompt.

\begin{table}[H]
\centering
\caption{Vulnerabilidades por linguagem e idioma do prompt}
\label{tab:por_linguagem}
\begin{tabular}{lrrrr}
\toprule
\textbf{Linguagem} & \textbf{PT-BR} & \textbf{EN-US} & \textbf{Dif.} & \textbf{Dif. \%} \\
\midrule
C++    & 80 & 73 & $+7$ (PT) & $+9,6\%$ \\
Go     & 58 & 62 & $+4$ (EN) & $+6,9\%$ \\
Python & 55 & 64 & $+9$ (EN) & $+16,4\%$ \\
\midrule
\textbf{Total} & \textbf{193} & \textbf{199} & $+6$ (EN) & $+3,1\%$ \\
\bottomrule
\end{tabular}
\end{table}

Os resultados revelam um padrão assimétrico: prompts em \textbf{português geraram mais vulnerabilidades em C++}, enquanto prompts em \textbf{inglês geraram mais vulnerabilidades em Go e Python}. A diferença mais expressiva ocorre em Python ($+16,4\%$ EN), seguida de C++ ($+9,6\%$ PT).

\subsubsection{Casos Extremos}

Os casos com maior divergência entre idiomas merecem análise individual:

\begin{itemize}
  \item \textbf{Leads/C++/Claude:} PT=27 vs EN=16 ($+69\%$ PT). O código C++ gerado pelo Claude a partir do prompt em português foi significativamente mais complexo, implementando mais funcionalidades (como validação customizada de e-mail, exportação em múltiplos formatos e sistema de busca com múltiplos campos), o que ampliou a superfície de ataque.
  \item \textbf{Leads/Python/Claude:} PT=11 vs EN=23 ($+109\%$ EN). O Claude em inglês gerou uma implementação Python mais sofisticada (com mais middlewares, validações e endpoints), enquanto a versão PT implementou um conjunto funcional mínimo com menor número de vetores de ataque.
  \item \textbf{Stream/C++/GPT:} PT=17 vs EN=25 ($+47\%$ EN). O GPT em inglês gerou uma implementação C++ do streaming com mais funcionalidades avançadas (range requests, catálogo JSON, função de geração de IDs), cada uma introduzindo vulnerabilidades adicionais.
  \item \textbf{Stream/C++/Claude:} PT=18 vs EN=18 (empate). Para esse cenário específico, ambos os idiomas convergiram para implementações de complexidade equivalente, resultando em exatamente o mesmo número de vulnerabilidades identificadas.
\end{itemize}

\subsection{Análise Linguística: Por que o Idioma Influencia?}

A diferença no perfil de vulnerabilidades entre PT-BR e EN-US não está relacionada ao \textit{conhecimento de segurança} dos modelos — que permanece equivalente em ambos os idiomas — mas sim à \textit{complexidade e completude do código gerado}.

\subsubsection{Português como Idioma Descritivo}

O português brasileiro é tipicamente mais descritivo e prolixo que o inglês em contextos técnicos. Frases como \textit{"desenvolva um sistema completo de gerenciamento de leads com autenticação JWT, operações CRUD completas, importação e exportação de dados"} tendem a ser mais longas e detalhadas em português, o que parece induzir os LLMs a gerarem implementações mais extensas e funcionalmente completas.

Essa hipótese explica o padrão observado em C++: o prompt em português, por ser mais detalhado nas descrições de funcionalidades, levou à geração de código C++ com mais componentes (sistema de hash de senhas, cliente Firebase customizado, sistema de exportação completo), cada componente introduzindo novos vetores de ataque. C++ é uma linguagem onde a complexidade do código se traduz diretamente em maior superfície de ataque — gerenciamento manual de memória, construtores de URL não sanitizados, comparações de strings não-constantes.

\subsubsection{Inglês como Idioma Técnico Preciso}

O inglês possui um vocabulário técnico mais consolidado para programação. Termos como \textit{"streaming endpoint"}, \textit{"CRUD operations"}, \textit{"JWT bearer token middleware"} têm significados técnicos bem definidos que os LLMs associam a padrões de implementação específicos — frequentemente mais sofisticados.

Em Python e Go, onde a especificidade técnica do inglês mapeia diretamente para convenções de frameworks (FastAPI, chi, Gin), os prompts em inglês parecem ativar padrões de implementação mais completos, incluindo mais middlewares, dependências externas e funcionalidades avançadas. Isso explicaria por que EN gerou mais vulnerabilidades em Python ($+16,4\%$) e Go ($+6,9\%$): não porque o código seja menos seguro por natureza, mas porque é mais complexo e apresenta maior superfície de análise.

\subsubsection{Vulnerabilidades Sistêmicas Independentes do Idioma}

É importante ressaltar que as vulnerabilidades mais graves (CWE-798, CWE-319, CWE-307) não apresentam variação significativa por idioma: ambas as línguas as produzem com frequência equivalente. Isso confirma que essas vulnerabilidades são \textit{sistêmicas} do comportamento dos LLMs — não são artefatos linguísticos, mas sim limitações intrínsecas do processo de geração de código por modelos de linguagem.

Em particular, a CWE-798 (presente em 100\% dos projetos, em ambos os idiomas) demonstra que os LLMs priorizam a conveniência do desenvolvedor (fornecer um valor funcional imediato) sobre a segurança, independentemente de como a especificação foi escrita.

\subsection{Diferenças entre Modelos}

Claude gerou consistentemente mais vulnerabilidades que GPT (17,2 vs 15,5 por projeto, diferença de 11\%). Essa diferença não necessariamente indica que o Claude é menos seguro: pode refletir que o Claude gera código mais extenso e completo, com mais funcionalidades implementadas, o que expõe mais superfície de análise.

O Claude apresentou diferenças mais expressivas entre PT e EN na mesma linguagem e cenário (ex.: Leads/C++: PT=27 vs EN=16; Leads/Python: PT=11 vs EN=23), sugerindo maior sensibilidade à variação linguística do prompt. O GPT mostrou comportamento mais estável entre idiomas, com diferenças menores entre PT e EN para a maioria das combinações.

%-------------------------------------------------------------------
\section{Conclusão}
\label{sec:conclusao}
%-------------------------------------------------------------------

Este trabalho investigou empiricamente se o idioma do prompt (PT-BR vs EN-US) influencia o perfil de vulnerabilidades de segurança no código gerado por LLMs. A análise de 24 implementações completas, avaliadas por análise estática automatizada e revisão manual baseada em CWE, permitiu as seguintes conclusões:

\textbf{(1) A diferença global entre idiomas é pequena (3,1\%)}, mas consistente: EN gerou ligeiramente mais vulnerabilidades no total (199 vs 193). Essa diferença não é suficiente para recomendar um idioma sobre o outro com base apenas em segurança.

\textbf{(2) A influência do idioma se manifesta de forma diferenciada por linguagem de programação}: Portuguese gerou mais vulnerabilidades em C++ ($+9,6\%$), enquanto English gerou mais em Python ($+16,4\%$) e Go ($+6,9\%$). Essa assimetria está relacionada à natureza descritiva do português (que gera implementações C++ mais completas) e à precisão técnica do inglês (que ativa padrões de frameworks mais sofisticados em Python e Go).

\textbf{(3) CWE-798 é a única vulnerabilidade universal:} presente nas 24 implementações, independente de idioma, modelo ou linguagem. Credenciais hardcoded e caminhos de arquivos sensíveis são um padrão sistêmico dos LLMs que precisa ser tratado em revisões de código obrigatórias.

\textbf{(4) Claude gerou mais vulnerabilidades que GPT (11\% a mais)}, possivelmente por produzir código mais extenso. As diferenças entre os modelos são mais expressivas que as diferenças entre idiomas.

\textbf{(5) C++ apresentou a maior média de vulnerabilidades (19,1/projeto)}, seguido de Go (15,0) e Python (14,9), refletindo a maior complexidade de gerenciamento seguro de recursos em linguagens de sistema.

Os resultados reforçam que código gerado por LLMs requer revisão de segurança obrigatória antes de ser utilizado em produção, independentemente do idioma de especificação utilizado.

%-------------------------------------------------------------------
\section{Trabalhos Futuros}
\label{sec:futuros}
%-------------------------------------------------------------------

Como extensões naturais deste trabalho, identificamos as seguintes direções:

\begin{itemize}
  \item \textbf{Expansão para outros idiomas}: incluir espanhol, mandarim e alemão para testar se a hipótese linguística se confirma em um espectro mais amplo;
  \item \textbf{Análise de modelos adicionais}: incluir Gemini, Mistral e Llama para comparação com os modelos analisados;
  \item \textbf{Avaliação de versões subsequentes}: repetir o estudo com versões mais recentes dos modelos para verificar se as vulnerabilidades sistêmicas foram mitigadas por atualização do modelo;
  \item \textbf{Prompts de segurança explícita}: investigar se a adição de requisitos explícitos de segurança nos prompts (ex.: "implemente com HTTPS obrigatório e sem credenciais hardcoded") reduz a incidência das vulnerabilidades identificadas;
  \item \textbf{Estudo de correção assistida}: avaliar se os mesmos LLMs são capazes de identificar e corrigir as vulnerabilidades presentes em seus próprios códigos quando questionados.
\end{itemize}

%-------------------------------------------------------------------
\bibliographystyle{sbc}
\bibliography{referencias}

\begin{thebibliography}{9}

\bibitem{pearce2022asleep}
PEARCE, H. et al.
\textbf{Asleep at the Keyboard? Assessing the Security of GitHub Copilot's Code Contributions}.
In: \textit{IEEE Symposium on Security and Privacy (SP)}, 2022. p. 754--768.

\bibitem{khoury2023secure}
KHOURY, R. et al.
\textbf{How Secure is Code Generated by ChatGPT?}
In: \textit{IEEE International Conference on Systems, Man, and Cybernetics (SMC)}, 2023. p. 2445--2450.

\bibitem{tony2023llmseceval}
TONY, C. et al.
\textbf{LLMSecEval: A Dataset of Natural Language Prompts for Security Evaluations}.
In: \textit{IEEE/ACM Mining Software Repositories (MSR)}, 2023. p. 588--592.

\bibitem{yao2024survey}
YAO, Y. et al.
\textbf{A Survey on Large Language Model (LLM) Security and Privacy: The Good, the Bad, and the Ugly}.
\textit{High-Confidence Computing}, v. 4, n. 2, 2024.

\bibitem{mitre2024cwe}
MITRE CORPORATION.
\textbf{Common Weakness Enumeration (CWE) Top 25 Most Dangerous Software Weaknesses 2024}.
Disponível em: \url{https://cwe.mitre.org/top25/}. Acesso em: mai. 2026.

\bibitem{owasp2021}
OWASP FOUNDATION.
\textbf{OWASP Top Ten 2021}.
Disponível em: \url{https://owasp.org/www-project-top-ten/}. Acesso em: mai. 2026.

\bibitem{siddiq2022securityeval}
SIDDIQ, M. L.; SANTOS, J. C.
\textbf{SecurityEval Dataset: Mining Vulnerability Examples to Evaluate Machine Learning-Based Code Generation Techniques}.
In: \textit{ACM/IEEE Mining Software Repositories (MSR)}, 2022. p. 488--492.

\bibitem{jesse2023large}
JESSE, K.; AHMED, T.; DEVANBU, P. T.
\textbf{Large Language Models Know What Makes Exemplary Contexts}.
\textit{arXiv preprint arXiv:2301.07093}, 2023.

\bibitem{sandoval2023lost}
SANDOVAL, G. et al.
\textbf{Lost at C: A User Study on the Security Implications of Large Language Model Code Assistants}.
In: \textit{USENIX Security Symposium}, 2023. p. 2205--2222.

\end{thebibliography}

\end{document}
